\subsection{Singleton pendants and actual smaller retracts}

We call $p$ a \emph{source singleton pendant with neighbor $q$}
if $p\ne q$, the hom-set $C(p,q)$ has exactly one element, and
every other strict hom-set incident to $p$ is empty. In particular
$C(q,p)=\varnothing$. A sink singleton pendant has the opposite
conditions. The corner at $p$ need not initially be trivial.
These definitions concern the whole category, not merely a chosen
subgraph of its support.

\begin{lemma}[Proved: integral recognition of a retracted subspace]
\label{CP2:A-retract-relative-products}
Let $i:Y\hookrightarrow X$ be a CW-subcomplex inclusion with a
continuous retraction $r:X\to Y$, and suppose $X\simeq\CP^2$.
If every product of positive-degree integral relative cohomology
classes in $H^*(X,Y;\mathbb Z)$ is zero, then $i$ is a homotopy
equivalence. The same conclusion holds if $i$ is a rational
homology equivalence.
\end{lemma}
\begin{proof}
The identity $ri=\operatorname{id}_Y$ makes each integral homology
group of $Y$ a direct summand of the corresponding group of $X$.
The space $Y$ is connected and simply connected, since it is the
image of a connected space under a retraction and its fundamental
group is a retract of the trivial group. Its only possible positive
integral homology groups are $0$ or $\mathbb Z$ in degrees two
and four. Universal coefficients give the same alternatives for
positive integral cohomology.

First assume the relative-product hypothesis. If $H^2(Y;\mathbb Z)=0$,
the generator $u\in H^2(X;\mathbb Z)$ lifts to a relative class.
Naturality of the relative cup product would give $u^2=0$, contrary
to the projective-plane ring. Thus $H^2(Y;\mathbb Z)=\mathbb Z$.
The split injection $r^*:H^2(Y;\mathbb Z)\to H^2(X;\mathbb Z)$
is consequently an isomorphism. Choose $v$ with $r^*v=u$.
Then $r^*(v^2)=u^2$, so the split injection in degree four is
also an isomorphism. Thus $i$ induces an isomorphism on all
integral homology groups. The homology form of Whitehead's
theorem for simply connected CW complexes, in the exact form
of \cite[Corollary~4.33]{CP2-Hatcher2002}, makes $i$ a homotopy
equivalence.

Under the second hypothesis, the same direct-summand description
and rational homology isomorphism already force the split maps
in degrees two and four to be integral isomorphisms. The same
Whitehead argument applies.
\end{proof}

\begin{proposition}[Proved: simplifying a rationally acyclic pendant corner]
\label{CP2:A-singleton-pendant-corner-retract}
Let $p$ be a source or sink singleton pendant of a finite category
$C$, with neighbor $q$. Put $L=C(p,p)$. If $BL$ is rationally
acyclic, replace $L$ by its identity and keep all other arrows
and all remaining composition values, obtaining a category $D$.
There are actual functors
\[
 i:D\hookrightarrow C,\qquad R:C\longrightarrow D,
 \qquad Ri=\operatorname{id}_D.
\]
The inclusion $Bi$ is a rational homology equivalence. If
$BC\simeq\CP^2$, then $Bi$ is a homotopy equivalence, and
$D$ has exactly $|\operatorname{Mor}C|-(|L|-1)$ arrows.
\end{proposition}
\begin{proof}
The functor $R$ fixes every object and every arrow outside $L$,
and sends every element of $L$ to $1_p$. This respects the
corner multiplication. It also respects all products involving
the sole cross arrow, since both corner actions on its singleton
hom-set are forced to be trivial. There are no other arrows
incident to $p$. Thus $D$ is a subcategory and the displayed
functors and their retraction identity are literal.

Write $E=C\setminus\{p\}$ for the full deletion and $N=C(q,q)$.
A nerve simplex containing $p$ cannot contain any strict object
other than $q$: a further object change would compose to an
additional strict arrow incident to $p$. Consequently the nerves
are unions of subcomplexes with intersection $NN$:
\[
 BC=BE\cup_{BN}B C|_{\{p,q\}},\qquad
 BD=BE\cup_{BN}B D|_{\{p,q\}}.
\]
The unique cross arrow identifies the first attached piece
naturally with $BL*BN$ in the source case, or $BN*BL$ in the
sink case. The second attached piece is the corresponding cone
on $BN$. The inclusion is the identity on $BE$ and $BN$, and
on the attached piece it is induced by the identity-vertex
inclusion $*\to BL$.

Because $BL$ is rationally acyclic, its join with the nonempty
connected space $BN$ is rationally acyclic, by the reduced join
homology formula over $\mathbb Q$. The cone is contractible.
The map between these connected attached pieces is therefore a
rational homology equivalence. The natural Mayer--Vietoris
sequences show that $BD\to BC$ is a rational homology
equivalence. If $BC\simeq\CP^2$, apply the second assertion of
Lemma~\ref{CP2:A-retract-relative-products}. Only the $|L|-1$
nonidentity arrows of $L$ were removed, proving the count.
\end{proof}

\begin{proposition}[Proved: folding a trivial singleton pendant]
\label{CP2:A-singleton-pendant-fold}
Let $p$ be a source singleton pendant of $C$, with trivial corner,
neighbor $q$, and sole cross arrow $f:p\to q$. Put $N=C(q,q)$
and $D=C\setminus\{p\}$. Suppose some object $r\in D$ admits
an arrow $g:r\to q$ satisfying
\[
 n\circ g=g\qquad(n\in N).
\]
Then an actual functor $R:C\to D$ fixes $D$, sends $p$ to $r$
and $f$ to $g$, and retracts the full inclusion $i:D\hookrightarrow C$.
If $BC\simeq\CP^2$, then $Bi$ is a homotopy equivalence and
$D$ has exactly two fewer arrows than $C$.

Dually, for a sink singleton pendant with sole arrow $f:q\to p$,
it suffices to have $g:q\to r$ with $g\circ n=g$ for every
$n\in N$. In particular, the hypothesis is automatic if an object
$r\ne p,q$ has a singleton hom-set to $q$ in the source case,
or from $q$ in the sink case. No condition on the corner or other
incidences of $r$ is required.
\end{proposition}
\begin{proof}
In the source case set $R(1_p)=1_r$ and use the stated assignments
on all other arrows. Products wholly within $D$ are unchanged.
The only additional nonidentity products involving $f$ are
$n\circ f=f$. They are preserved precisely by $n\circ g=g$.
The identity products at $p$ are also preserved. Thus $R$ is a
functor and $Ri=\operatorname{id}_D$.

As in the preceding proof, every simplex not wholly in $D$ lies
in the two-object piece on $p,q$. Since the corner at $p$ is
trivial, that piece is a cone on $BN$, attached along its base.
Consequently the quotient of the actual nerve pair is
\[
 BC/BD\simeq \operatorname{Cone}(BN)/BN\simeq\Sigma BN.
\]
The relative cohomology ring has zero products of positive-degree
classes. Lemma~\ref{CP2:A-retract-relative-products} therefore
gives the claimed homotopy equivalence. The deletion removes
only $1_p$ and $f$. Taking opposite categories proves the sink
case, with the stated reversed composition law. On a singleton
hom-set every endomorphism action is forced to fix its sole
element, proving the last assertion.
\end{proof}

\begin{corollary}[Proved: two cardinality obstructions at eighteen arrows]
\label{CP2:A-singleton-pendant-eighteen}
An eighteen-arrow category of type $\CP^2$ cannot have either:
\begin{enumerate}
\item a singleton pendant whose corner has order two, three or four;
\item a trivial singleton pendant and another singleton hom-set
in the same direction at its neighbor, with a different other object.
\end{enumerate}
\end{corollary}
\begin{proof}
All corner monoids of orders at most four have rationally acyclic
classifying spaces by Corollary~\ref{CP2:A-five-types}.
In the first case the corner replacement loses at least one
arrow, and in the second the deletion loses two. Each construction
would give a category of type $\CP^2$ with at most seventeen
arrows, contrary to Theorem~\ref{CP2:A-global-lower-eighteen}.
\end{proof}
